% Folder 144, batch 04 (archivists' pages 61-80).
% Pass: Opus 5.5 (claude-opus-5-5), 2026-10-03 — first pass, unchecked against the pages by a human.
\documentclass[11pt,a4paper]{article}
\input{../preamble/grothendieck}

\folder{144}
\batch{4}
\pages{61}{80}
\foldertitle{[Suite autour de Teichmüller dont] Teichmülleries, 1981-1982 : notes manuscrites (1981-1983, s.d.), lettres (1981, s.d.).}
\dating{1981-1983}
\watermark{Édition de démonstration}

\begin{document}

\page{61}\note{p. 6 de l'auteur. La page reprend une phrase dont le début est sur la page précédente, hors de ce lot : la comparaison entre cartes et revêtements ramifiés de $\mathbb{P}^1(\mathbb{C})$ est commencée avant la p. 61.}
C'est ainsi que \add{la cat.\ des} les cartes cellulaires triangulées pondérées
apparaissent comme équivalentes \add{à} : celles des \add{surjectifs}
revêtements \add{ramifiés} de \struck{\ill{}} $\mathbb{P}^1(\mathbb{C})$, ramifiés
au plus en les points $0, 1, \infty$. Ces revêtements sont de façon
\uncertain{connue} \ill{} des courbes algébriques (propres et lisses) sur
$\mathbb{C}$, et la \struck{carte} \add{donnée par $X$} est l'image inverse de
\add{l'équateur} $\mathbb{P}^1(\mathbb{R})$, les sommets de type $i$ \add{formant}
les fibres en $X$ au-dessus du point $i \in \{0, 1, \infty\} \subset
\mathbb{P}^1(\mathbb{C})$, les ordres \uncertain{réduits} les pts de ramification ;
les arêtes sont les composantes connexes des images inverses des segments ouverts
$]1, \infty[$, $]\infty, 0[$, $]0, 1[$. \add{(subdivisions barycentriques des)}

Le cas des cartes ordinaires est celui où \ill{} :
\[
  \rho_1^2 = 1 \;) ,
\]
cette dernière se récupère en termes du revêtement en prenant l'image inverse du
segment $[0, 1] \subset U_{0,3} \subset \mathbb{C}$, les sommets
(\uncertain{ordinaires}) sont les pts au-dessus de $0$, \uncertain{tandis} que les
points au-dessus de $1$ (comp.\ $\infty$) sont par définition les centres des
arêtes resp.\ centres des faces. On trouve la carte \add{orientée}
« universelle » \struck{orientable} sur $\mathbb{P}^1(\mathbb{C})$, comme
\struck{le segment},
\[
  \text{(21)} \qquad S_0 = \{0\} \subset K_\infty = [0, 1]
\]
\note{devant (21), un numéro biffé illisible.}

\page{62}
et la subdivision barycentrique, donnée par
\[
  \text{(22)} \qquad
  \left\{
  \begin{array}{l}
    S_1 = \{1\}, \quad S_\infty = \{\infty\}, \quad
    K = \mathbb{P}^1(\mathbb{R}) = K_0 \cup K_1 \cup K_\infty \\
    \text{où} \quad K_0 = [1, \infty], \quad K_1 = [\infty, 0], \quad
    \underbrace{K_\infty = [0, 1]}_{\text{comme avant, cf.\ (21)}}
  \end{array}
  \right.
\]
\drawing{Une sphère, coupée par l'équateur ; sur l'équateur, les points $0$, $1$, $\infty$ et les arcs $K_0$, $K_1$, $K_\infty$ (ce dernier, de $0$ à $1$, renforcé) ; l'hémisphère supérieur marqué $\Gamma_+$, l'inférieur $\Gamma_-$, avec une flèche d'orientation.}

pour le cas de la carte orientée triangulée pondérée universelle, avec $K_0, K_1,
K_\infty$ comme arêtes de type $0, 1, \infty$ respectivement. Les deux
\struck{faces} \add{triangles-repères} dans la situation universelle orientée
triangulée pondérée sont l'hémisphère nord $\Gamma_+$ et l'hémisphère sud
$\Gamma_-$, qui apparaissent comme des « simplexes » topologiques, ayant comme
sommets $0, 1, \infty$, et comme arêtes $K_0, K_1, K_\infty$ opposées à ceux-ci.
Quant aux faces de la situation initiale (21), il y a une seule face, découpée de
la sphère $\mathbb{P}^1(\mathbb{C})$ suivant le segment $K_\infty = [0, 1]$, avec
\add{comme bord l'}\struck{une} arête \uncertain{unique}
$0'\,(\tfrac12)_-\,1'\,(\tfrac12)_+\,0'$ \struck{\ill{}} \add{parcourue}
une boucle fermée \uncertain{ayant} comme seul sommet $0' \mapsto 0$, comme
« milieu » le point $1' \mapsto 1$, les deux arcs joignant $0'$ à $1'$ sur le bord
de la découpure étant ceux qui passent par $\Gamma_-$ (contenant
$(\tfrac12)_-$) et $\Gamma_+$ respectivement. Vu depuis $\infty$, centre de
l'unique face, en ici \uncertain{collant} les deux bords de la découpure,
\add{celle-ci} \struck{\ill{}} apparaît donc comme un disque de centre $\infty$,
divisé par les deux rayons \struck{$(0', \infty)$} $(\infty, 0')$ et
$(\infty, 1')$, joignant $\infty$ à $0 \sim 0'$ et à $1 \sim 1'$
respectivement, en les deux repères $\Gamma_+$ et $\Gamma_-$ de la subdivision
barycentrique, $\Gamma_+$ ayant \struck{\ill{}} comme orientation propre celle
induite par $f$, \struck{\ill{}} et $\Gamma^-$ ayant comme orientation propre
l'orientation opposée de celle induite par $f$. \uncertain{Par la}

\drawing{Dans la marge gauche, deux figures. En haut, la sphère avec la découpure le long de $[0,1]$ : les bords $0'$, $1'$, $(\tfrac12)_+$, $(\tfrac12)'_-$, et des flèches parcourant le bord. En bas, le disque $f$ vu depuis son centre $\infty$ : diamètre horizontal pointillé de $1'$ (à gauche) à $0'$ (à droite), $\Gamma^+$ au-dessus, $\Gamma^-$ au-dessous, avec leurs flèches d'orientation ; sur le cercle, $(\tfrac12)'_+$ en haut et $(\tfrac12)'_-$ en bas, arcs marqués $K'^{+}_{\infty}$ et $K'^{-}_{\infty}$ ; une accolade $f$ à gauche.}

\marginal{L'orientation marquée sur le bord de la découpure est celle induite par
l'unique face $f$, \uncertain{non} l'orientation habituelle héritée de
$\mathbb{P}^1(\mathbb{C})$.}

\marginal{Attention, si on place $\Gamma^+$ au-dessus dans l'« horizontale »
$0'\infty 1'$, $0'$ doit être placé à droite \uncertain{et} $1'$ à gauche !}

\page{63}\note{p. 7 de l'auteur.}
théorie de la représentation conforme, ce disque holomorphe doit être isomorphe
au disque standard $\mathbb{D}$, avec un isomorphisme unique
\[
  g : \mathbb{D} \xrightarrow{\ \sim\ } f \qquad
  \begin{array}{l}
    0 \longmapsto \infty \quad (\text{centre de } f) \\
    1 \longmapsto 0' \quad (\text{unique sommet de } f)
  \end{array}
\]
et tel que de plus
\[
  \left\{
  \begin{array}{l}
    g(-1) = 1' \\
    g(i) = (\tfrac12)'_+ \\
    g(-i) = (\tfrac12)'_-
  \end{array}
  \right.
\]
\drawing{Dans la marge, le disque unité : les points $1$ (renforcé), $-1$, $i$, $-i$ et le centre $0$, une flèche d'orientation sur le cercle.}

Comment trouver $g$ ? On \struck{trouve} \add{voit} que le disque $\mathbb{D}$
est l'hémisphère positif de la « sphère » standard \struck{$X_1$}
\add{$= \mathbb{P}^1(\mathbb{C})$}, qui \struck{\ill{}} a comme seul sommet $1$,
comme seule arête l'équateur $\mathbb{P}^1(\mathbb{R})$, comme centres des
\add{deux} faces $0$ et $\infty$ (qui forment l'axe polaire), comme milieu
d'arête $-1$. Alors \add{$X_1$} \struck{\ill{}} est un revêtement d'ordre $2$ de
$X_0 = \mathbb{P}^1(\mathbb{C})$, ramifié \struck{en dehors} au plus en
$\{0, 1, \infty\}$ (l'ordre du revêtement est le nb de repères
\add{directs}, qui est $2$)

\drawing{Dans la marge, une sphère : $0$ au pôle nord, $\{\infty\}$ au pôle sud, $1$ sur l'équateur, une flèche sur l'équateur.}

où \struck{\ill{}}
\[
  g : \quad
  \begin{array}{l}
    1 \longmapsto 0 \\
    -1 \longmapsto 1 \\
    0, \infty \longmapsto \infty
  \end{array}
\]
ce qui \uncertain{épuise} les \uncertain{fibres} en $0, 1, \infty$, et montre que
la ramification est $2$ sur $0, 1$, et pas de ramification \struck{\ill{}}
\add{p.\ ex.} \ill{} la situation \add{analogue,}
$g' : \mathbb{P}^1(\mathbb{C}) \to \mathbb{P}^1(\mathbb{C})$, \uncertain{mais}
où
\[
  g' : \quad
  \begin{array}{l}
    0 \longmapsto 0 \\
    \infty \longmapsto \infty \\
    1, -1 \longmapsto 1
  \end{array}
\]
on a \struck{\ill{}} évidemment $g'(z) = z^2$,

\page{64}\note{p. 8 de l'auteur.}
et on a
\[
  g(z) = \varphi\, g'(\psi(z))
\]
où $\varphi, \psi$ sont les transformations homographiques
\[
  \psi : \left\{
  \begin{array}{l} 1 \\ -1 \\ 0, \infty \end{array}
  \right.
  \longrightarrow
  \left\{
  \begin{array}{l} 0 \\ \infty \\ 1, -1 \end{array}
  \right.
  \qquad
  \varphi : \left\{
  \begin{array}{ll} 0 & 0 \\ \infty & 1 \\ 1 & \infty \end{array}
  \right.
\]
donc
\[
  \psi(z) = \frac{z-1}{z+1} \qquad
  \varphi(z) = \sigma_0(z) = \frac{z}{z-1}
\]
d'où
\[
  g(z) = \frac{\left(\frac{z-1}{z+1}\right)^2}
              {\left(\frac{z-1}{z+1}\right)^2 + 1}
       = \frac{(z-1)^2}{(z-1)^2 - (z+1)^2}
       = \frac{(z-1)^2}{-4z}
\]
\note{le dénominateur du premier membre se lit $+1$ ; la valeur de $\varphi$ et
le membre suivant demandent $-1$.}
\[
  \text{(23)} \qquad g(z) = -\frac{1}{4z}\,(z-1)^2
\]
réalise l'isomorphisme conforme du disque unité $\mathbb{D}$ sur le découpage
$f$ de $\mathbb{P}^1(\mathbb{C})$ suivant $[0, 1]$, avec les valeurs annoncées
plus haut…

\note{un trait horizontal sépare ici deux développements.}

Le groupe $\Pi_{0,3}^\tau$ apparaît comme un groupe fondamental mixte
\[
  \text{(24)} \qquad \Pi_{0,3}^\tau \simeq
  \Pi_1(U_{0,3}, \{1, \tau\}; P_+)
\]
\marginal{\textbf{NB} Ici $\rho_i = \lambda_{i+1} \lambda_{i-1}^{-1}$}
Introduisant les chemins \add{$\lambda_i$ ($i \in \{0, 1, \infty\}$)} de
$P_-$ à $P_+$, formés des demi-arcs de grand cercle passant par les $Q_i$
(\add{$Q_i$,} milieux des arcs de grand cercle joignant $j$ à $k$, où
$\{i, j, k\} = \{0, 1, \infty\}$), on arrive, en tant que flèches de
$\Pi_1(U_{0,3}, \{1, \tau\})$ :
\[
  \tau_i = (\underbrace{\lambda_i}_{\tau P_+ = P_- \to P_+},
  \underbrace{\tau}_{\{1, \tau\}})
\]
\drawing{(24) Dans la marge, une sphère : $P_+$ au pôle nord ; trois demi-méridiens $\lambda_0$, $\lambda_1$, $\lambda_\infty$ qui y aboutissent, passant par $Q_0 = 2$, $Q_1 = -1$, $Q_\infty = \tfrac12$ ; les points $0$, $1$, $\infty$ sur l'équateur ; des parties des arcs en pointillé.}
\note{le $Q_0 = 2$ de la figure se lit ainsi ; les deux annotations sous $\lambda_i$ et $\tau$ sont transcrites comme accolades.}

\page{65}
Cette interprétation montre que la catégorie des cartes \add{isotropiques}
triangulées pondérées (de type $\{0, 1, \infty\}$) est équivalente :
\add{une ss-cat.\ pleine de} celle des \add{$\{1, \tau\}$-}revêtements
\add{finis} étales de $U_{0,3}$, \struck{tels que} \add{savoir ceux}
satisfaisant à une condition supplémentaire (savoir que sur l'ensemble opposé
correspondant, $\tau_0, \tau_1, \tau_\infty$ opèrent sans pts fixes). Mais on
voit \uncertain{aisément}, \struck{\ill{}} que pour le \add{$\{1, \tau\}$}
revêtement $U'$ de $U_{0,3}$, les points fixes de $\tau_i$ correspondent aux
composantes connexes de $U'^\tau$ qui sont au-dessus de l'arc de
$U_{0,3}(\mathbb{R})$ opposé à $P_i$ i.e.\ \struck{autour} de centre $Q_i$ ---
donc la condition signifie que \struck{\ill{}} l'on a
\[
  \text{(25)} \qquad U'^\tau = \emptyset \quad \text{i.e.} \quad
  X'^\tau = \emptyset
\]
où $X'$ est le revêtement ramifié de $\mathbb{P}^1(\mathbb{C})$ qui prolonge
$U'$. [\textbf{NB} \struck{\ill{}} a priori $X'^\tau$ est une courbe analytique
réelle \struck{plongée} \add{contenue} dans $X'$, qui \struck{\ill{}}
\add{donc} est vide dès qu'elle se trouve dans
$U' = X' \smallsetminus \text{ramif.}$ \struck{(\ill{})} l'est.)]
\add{NB.} (En termes de la carte pondérée initiale, $X'$ est le revêtement de
ses orientations \add{locales}, $\tau_{X'}$ \struck{est} l'automorphisme défini
par \ill{} : l'orientation $\mapsto$ l'orientation opposée…, qui en effet doit
être sans points fixes).

Indépendamment de la condition restrictive (25), la donnée du $\{1, \tau\}$
revêtement \add{ramifié} de $\Sigma = \mathbb{P}^1(\mathbb{C})$,

\page{66}\note{p. 9 de l'auteur.}
étale au-dessus de \add{$\Sigma^* =$} $U_{0,3}$, \struck{revient}
\add{i.e.} la donnée d'un revêtement ramifié de $\Sigma$ étale au-dessus de
$U_{0,3}$, plus une op.\ de $\{1, \tau\}$ dessus compatible avec l'opération
sur $\Sigma$, revient : la donnée d'une \emph{donnée de descente} de
$\mathbb{C}$ à $\mathbb{R}$, sur un revêtement ramifié, étale au-dessus de
$U_{0,3\,\mathbb{C}}$, du \add{$\mathbb{C}$-}schéma $\Sigma_{\mathbb{C}}$. Donc
la catégorie des revêtements mixtes \add{est} équivalente : celle des
revêtements ramifiés du schéma \struck{$U_{0,3}$}
\add{$\Sigma_{\mathbb{R}} = \mathbb{P}^1_{\mathbb{R}}$}, étales au-dessus de
\[
  \text{\struck{$U$}} \ U_{0,3\,\mathbb{R}} = \Sigma^*_{\mathbb{R}}
  = \Sigma_{\mathbb{R}} \smallsetminus \{0, 1, \infty\} .
\]
On trouve ainsi un \emph{foncteur pleinement fidèle}
\[
  \text{(26)} \quad
  \begin{array}{c}
    \text{cartes triangulées pondérées \textit{finies}} \\
    \text{de type } \{0, 1, \infty\}
  \end{array}
  \hookrightarrow
  \begin{array}{l}
    \text{revêtements étales \textit{finis}} \\
    \text{du schéma } \mathbb{P}^1_{\mathbb{R}} \smallsetminus \{0,1,\infty\} \\
    (\simeq \text{rev.\ ramifiés de } \mathbb{P}^1_{\mathbb{R}}, \\
    \ \text{étales au-dessus de} \\
    \ U^*_{03\,\mathbb{R}} = \mathbb{P}^1_{\mathbb{R}} \smallsetminus \{0,1,\infty\})
  \end{array}
\]
\note{dans (26), « finies » et « finis » sont ajoutés au-dessus de la ligne.}

L'image essentielle est formée des rev.\ $X'$ tels que \struck{\ill{}}
\[
  \text{(27)} \qquad X'(\mathbb{R}) = \emptyset
\]
interprétation qui en \uncertain{fait} \uncertain{traduit} (25) puisque
\[
  X'(\mathbb{R}) = \underbrace{X'(\mathbb{C})^\tau}_{\text{\struck{\ill{}} $X'$ dans (25)}} .
\]

Le foncteur en sens inverse est donné en associant : à un rev.\
\add{ramifié} $X'$ de $X$, \uncertain{étale sur} $U_{0,3}$, \uncertain{la carte}
$X'(\mathbb{C})$, munie de l'image \uncertain{inverse}
$X'(\mathbb{C}) \,|\, (\mathbb{P}^1_{\mathbb{R}} = \Sigma_{\mathbb{R}})$ de
l'ens.\ des pts réels de $\mathbb{P}^1(\mathbb{C}) = \Sigma$,

\page{67}
\marginal{\struck{(B4.\ B$_0$4)}}\note{cette mention biffée, en haut à gauche, est
à l'encre bleue ; la lecture en est douteuse.}
qui définit une triangulation (orientée) avec partition par les images inverses
des pts $0, 1, \infty$ \struck{par définition} (qui forment \add{une partition
de l'ens.\ des} sommets de la triangulation, défini comme l'ensemble des pts
au-dessus de $\{0, 1, \infty\}$), et en passant au quotient par la conjugaison
complexe --- ce qui donne une surface \uncertain{sans} bord, grâce à la
condition (27). Dans le cas où on ne passe pas par cette condition, l'espace
quotient
\[
  \text{(28)} \qquad X = X'(\mathbb{C}) / \tau
\]
est une surface à bord compacte (et munie d'une structure conforme…), et on
pourrait étendre \add{(de façon essent.\ unique)} la définition des cartes
\add{finies} triangulées pondérées aux surfaces compactes à bord, de telle façon
que (26) devienne une équivalence de catégories.

(4) On peut introduire, sur une carte (orientée ou non), des conditions de la
forme
\[
  \text{(29)} \qquad \rho_i^{\nu_i} = 1
  \qquad \text{\struck{type}} \ \text{où} \quad
  \nu_i \in \overline{\mathbb{N}^*},
\]
qui signifie donc que

\page{68}\note{p. 10 de l'auteur.}
dans la carte triangulée pondérée envisagée, les sommets de type $i$
\struck{\ill{}} \add{ont un} ordre réduit \emph{divisant} $\nu_i$.
Une telle condition peut être introduite pour un ou plusieurs indices $i$. Dans le
cas des cartes ordinaires (interprétées comme définissant une triangulation
pondérée, via la subdivision barycentrique) une relation
\[
  \text{(30)} \qquad \rho_0^{\nu_0} = 1 \qquad
  (\text{resp.}\ \rho_\infty^{\nu_\infty} = 1)
\]
signifie que les sommets \struck{\ill{}} ont d'ordre divisant $\nu_0$
(resp.\ que les faces sont d'ordre divisant $\nu_\infty$). Quant à la relation
\[
  \text{(31)} \qquad \rho_1^{\nu_1} = 1 \qquad )
\]
comme \ill{} déjà $\rho_1^2 = 1$, ou bien $\nu_1$ \add{(\uncertain{tel} \ill{})}
et cette relation est automatiquement satisfaite, ou bien $\nu_1$ est impair, et
elle implique
\[
  \rho_1 = 1
\]
ce qui signifie que la carte correspond à un revêtement de $\mathbb{P}^1_{\mathbb{C}}$
ramifié seulement en $0$ et $\infty$, donc : un revêtement « cyclotomique »
$z \mapsto z^n$ ($n \in \mathbb{N}$), \struck{cas qu} donnant lieu : la carte
\uncertain{régulière} \add{sphérique} \uncertain{de groupe}
$\mathbb{Z}/n\mathbb{Z} \simeq \mu_n(\mathbb{C})$ ayant un seul sommet, une seule
face, et $n$ arêtes disposées circulairement (cas \ill{} isomorphe : se
\uncertain{démêle}, visiblement) --- cas qui est relativement trivial. Donc on
pourra
\drawing{En bas à gauche, une sphère : un sommet unique au pôle nord, d'où partent plusieurs méridiens, certains en pointillé (la carte cyclotomique à un sommet, une face et $n$ arêtes).}

\page{69}
\struck{\ill{}}
\note{Cette page porte un renvoi (deux traits obliques) vers un ajout encadré.}
que dans le cas des cartes ordinaires, les relations considérées sont du type (30),
\uncertain{l'exclusion de (31)}. [On considère que $\rho_i^\infty = 1$ par
définition, de sorte que \add{l'on} \struck{les cartes} \ill{}]
l'on peut considérer un système de conditions (29), où pour tout
$i \in \{0, 1, \infty\}$, on a
\[
  \text{(34)} \qquad \nu_i \in \overline{\mathbb{N}^*} = \mathbb{N}^* \cup \{+\infty\} .
\]
\note{le numéro se lit (34) ; il fait double emploi avec le (34) de la p. 70.}
Les cartes (finies) pondérées triangulées correspondantes sont déterminées par
l'action du groupe quotient
\[
  \text{(32)} \qquad \Pi_{0,3} / \langle \rho_0^{\nu_0}, \rho_1^{\nu_1},
  \rho_\infty^{\nu_\infty} \rangle \qquad (\text{cas orienté})
\]
ou par le quotient
\[
  \text{(32 bis)} \qquad \Pi_{0,3}^\tau / \langle \rho_0^{\nu_0}, \rho_1^{\nu_1},
  \rho_\infty^{\nu_\infty} \rangle \qquad (\text{cas non orienté})
\]
qui s'identifie : une extension de $\{1, \tau\}$ par (32). Ce dernier groupe
s'interprète comme un groupe fondamental « avec singularités »
\[
  \text{(33)} \qquad \Pi_{0,3} / \langle (\rho_i^{\nu_i}) \rangle \simeq
  \Pi_1(\mathbb{P}^1(\mathbb{C}), \nu_*; P_+)
\]
\[
  \nu_* = (\nu_0 R_0 + \nu_1 R_1 + \nu_\infty R_\infty) \qquad
  \text{singularités sur } \mathbb{P}^1_{\mathbb{C}}
\]
ou encore comme un groupe fondamental de la multiplicité (topologique, analytique
ou schématique, au choix --- mais dans ce dernier cas il s'agit du complété
profini…)

\page{70}\note{p. 11 de l'auteur.}
\[
  \text{(34)} \qquad \mathcal{X}(\mathbb{P}^1_{\mathbb{C}}, \nu_*) =
  \text{\struck{catégorie}}\ \text{topos des \struck{espaces} \add{surfaces}
  topologiques}
\]
\begin{quote}
\struck{(2 \ill{})} (\uncertain{non nécess.\ séparées}) ramifiées sur
$\mathbb{P}^1_{\mathbb{C}}$, avec indices de ramification, \struck{aux}
\add{en} les points au-dessus d'un $R_i$ ($i \in S$ = \uncertain{supp.}\
$\nu_*$), divisant $\nu(R_i) = \nu_i$
\end{quote}
\marginal{NB ce topos est une multiplicité topologique, loc.\ surface topologique,
elle est canoniquement orientée.}

i.e.\ on a un isom.\ canonique
\[
  \text{(35)} \qquad \Pi_{03} / \langle (\rho_i^{\nu_i}) \rangle \simeq
  \Pi_1(\mathbb{P}^1(\mathbb{C}), \nu_*; P_+) \overset{\text{déf}}{=}
  \Pi_1(\mathcal{X}(\mathbb{P}^1(\mathbb{C}), \nu_*), P_+) ,
\]
et de même
\[
  \begin{aligned}
  \text{(36)} \qquad \Pi_{0,3}^\tau / \langle (\rho_i^{\nu_i}) \rangle
  &\simeq \Pi_1(\mathbb{P}^1(\mathbb{C}), \{1, \tau\}, \nu_*; P_+) \\
  &\overset{\text{déf}}{=}
  \Pi_1(\mathcal{X}(\mathbb{P}^1(\mathbb{C}), \nu_*, \{1, \tau\}), P_+)
  \end{aligned}
\]
où
\[
  \text{(37)} \qquad \mathcal{X}(\mathbb{P}^1(\mathbb{C}), \nu_*, \{1, \tau\})
  = (\mathcal{X}(\mathbb{P}^1(\mathbb{C}), \nu_*), \{1, \tau\})
\]
topos mixte défini par le topos (34) avec l'opération naturelle de $\{1, \tau\}$
dessus.
\marginal{NB c'est encore une multiplicité topologique, localement une surface,
mais elle \uncertain{n'est} multiplicité n'est pas orientable…}

Passant aux complétés profinis, on trouve
\[
  \begin{aligned}
  \text{(38)} \qquad \widehat{\Pi}_{0,3} / \langle (\rho_i^{\nu_i}) \rangle
  &\simeq \Pi_1(\underbrace{\mathbb{P}^1_{\mathbb{C}}}_{\text{schéma}}, \nu_*; P_+) \\
  &\simeq \Pi_1(\mathbb{P}^1_{\overline{\mathbb{Q}}}, \nu_*; P_+)
  \end{aligned}
\]
\[
  \begin{aligned}
  \text{(39)} \qquad \widehat{\Pi}_{0,3}^\tau / \langle (\rho_i^{\nu_i}) \rangle
  &\simeq \Pi_1(\underbrace{\mathbb{P}^1_{\mathbb{R}}}_{\text{schéma}}, \nu_*; P_+) \\
  &\simeq \Pi_1(\mathbb{P}^1_{\overline{\mathbb{Q}} \cap \mathbb{R}}, \nu_*; P_+)
  \end{aligned}
\]
(et \struck{\ill{}} l'interprétation \add{(38)} montre qu'il y a une

\page{71}
opération extérieure de $\mathrm{Gal}(\overline{\mathbb{Q}}/\mathbb{Q})$ sur le
premier membre de (38), et en particulier sur $\widehat{\Pi}_{0,3}$ lui-même…).

Les cas les plus intéressants sont les \struck{deux} suivants :

I) $\nu_0 = \nu_1 = \nu_\infty = +\infty$ : on retrouve $\Pi_{0,3}$ et
$\Pi_{0,3}^\tau$, classifiant les cartes triangulées pondérées (resp.\ orientées).

II) $\nu_0 = \nu_\infty = +\infty$, $\nu_1 = 2$ : on retrouve les groupes
\[
  \Pi_{0,3}^\tau / \langle \rho_1^2 \rangle \quad \text{et} \quad
  \Pi_{0,3} / \langle \rho_1^2 \rangle ,
\]
classifiant les cartes ordinaires.

III) $\nu_0 = \infty$, $\nu_1 = 2$, $\nu_\infty = 3$, donnant les groupes
\[
  \Pi_{0,3}^\tau / \langle \rho_1^2, \rho_\infty^3 \rangle \quad \text{et} \quad
  \Pi_{0,3} / \langle \rho_1^2, \rho_\infty^3 \rangle ,
\]
classifiant les cartes ordinaires \emph{triangulaires} (\emph{pas} pondérées !),
i.e.\ dont toutes les faces sont \struck{des \uncertain{triangles}}
\add{i.e.\ sont des triangles} --- dans le cas non orienté, et les
\add{les} \ill{} \uncertain{orientées}. On trouve des isomorphismes
(via (35), (36))
\[
  \underbrace{\Pi_{0,3}^\tau / \langle \rho_1^2, \rho_\infty^3 \rangle}
    _{\text{groupes cartographiques}}
  \xleftarrow{\ \sim\ }
  \underbrace{\Pi_1(U_{0,3}, \mathfrak{S}_3 \times \{1, \tau\}; P_+)}
    _{\Pi_{0,3}^{\mathbb{D}_3, \tau}}
\]
\[
  \underset{\text{triangulés, non or.\ resp.\ orientés}}
    {\Pi_{0,3} / \langle \rho_1^2, \rho_\infty^3 \rangle}
  \xleftrightarrow{\ \sim\ }
  \overbrace{\Pi_1(U_{0,3}, \mathfrak{S}_3; P_+)}^{\Pi_{0,3}^{\mathbb{D}_3}}
\]
\note{les deux membres de gauche sont réunis par une accolade, sous laquelle
« groupes cartographiques triangulés, non or.\ resp.\ orientés » ; le second
membre de gauche est rattaché par un trait double à
$\{\rho_1, \rho_\infty \mid \rho_1^2 = \rho_\infty^3 = 1\}$, écrit dans la
marge.}
\marginal{Ces isomorphismes ne voilent pas une \uncertain{diversité}
d'opérations \ill{} \ill{} \ill{} \ill{} $\Pi_1(X, G; P)$ groupe fondamental
« \uncertain{orbifold} » \ill{} des \ill{} \uncertain{triangulaires}
\ill{} \ill{} \ill{} (40)}
\note{note marginale écrite en diagonale, en partie illisible ; elle renvoie
semble-t-il à un (40) à venir.}

\page{72}\note{p. 12 de l'auteur.}
dont le deuxième est donné par
\[
  \text{(41)} \qquad
  \boxed{\rho \longmapsto \rho_\infty^{-1}, \quad \sigma_\infty \longmapsto \rho_1}
  , \quad
  \underset{\sigma_\infty \rho}{\varepsilon_0} \longmapsto \rho_0
\]
\note{sous $\rho_\infty^{-1}$, une flèche et l'annotation « attention : l'exposant ! » ; sous $\varepsilon_0$, un signe $=$ vertical renvoyant à $\sigma_\infty\rho$.}
et le premier par
\[
  \text{(42)} \qquad
  \boxed{\tau_\infty \longmapsto \tau_\infty}, \quad
  \underset{\rho(\tau_\infty)}{\tau_0} \longmapsto
  \overbrace{\underbrace{\rho_\infty^{-1} \tau_\infty \rho_\infty}
    _{(\rho_\infty^{-1}\rho_1\rho_\infty^{-1}\rho_1)\tau_\infty}}
    ^{\tau_0\tau_1\tau_\infty\tau_1\tau_0} ,
\]
\[
  \underset{\rho^2(\tau_\infty)}{\tau_1} \longmapsto
  \overbrace{\rho_\infty^{-2} \tau_\infty \rho_\infty^{2}}
    ^{\tau_1\tau_0\tau_\infty\tau_0\tau_1}
  = \rho_\infty \tau_\infty \rho_\infty^{-1}
  \quad (\text{car } \rho_\infty^3 = 1 \text{ dans }
  \Pi_{03}^\tau / \langle \cdots \rangle)
\]
\note{au-dessus de la dernière accolade, entre parenthèses :
$= (\rho_\infty\rho_1\rho_\infty\rho_1)\tau_\infty$.}
\marginal{NB on utilise $\rho_1^2 = \rho_\infty^3 = 1$ et
$\tau_\infty(\rho_\infty) = \rho_0\rho_1 = \rho_1^{-1}\rho_\infty^{-1}\rho_1
= \rho_1\rho_\infty\rho_1$}
où
\[
  \rho, \sigma_i, \varepsilon_i \ (i \in \{0, 1, \infty\}) \quad \text{dans} \quad
  \Pi_1(U_{0,3}, \mathfrak{S}_3; P_+) = \Pi_{03}^{\mathbb{D}_3}
\]
sont définis de façon standard (on y reviendra dans un cas plus général un peu
plus bas), et
\[
  \tau_i \ (i \in \{0, 1, \infty\}) \in
  \underbrace{\Pi_1(U_{0,3}, \{1, \tau\}, P_+)}_{\Pi_{0,3}^\tau}
  \subset
  \underbrace{\Pi_1(U_{0,3}, \mathfrak{S}_3 \times \{1, \tau\}; P_+)}
    _{\Pi_{0,3}^{\mathbb{D}_3, \tau}}
\]
sont définis dans (24) (cf.\ \struck{\ill{}} aussi les générateurs de
$\Pi_{0,3}^\tau \simeq \mathfrak{S}\mathfrak{l}_2$ \uncertain{qui figurent} dans
(2)). Je vais revenir sur ces isomorphismes (40) dans un contexte un peu plus
général, ci-dessous.

Enfin, dans le développement précédent, remplaçant $U_{0,3}$ et l'action de
$\mathfrak{S}_3$ resp.\ $\mathfrak{S}_3 \times \{1, \tau\}$ dessus, par
\struck{celle de} le demi-plan de Poincaré $\Sigma^+$ et l'action de
$\mathrm{Sl}(2, \mathbb{Z})/\pm 1$ \add{\struck{Gl}} resp.\
$\mathrm{Gl}(2, \mathbb{Z})/\pm 1$ dessus (par l'action homographique
resp.\ homographique et « antihomographique »), on trouve aussi des
isomorphismes remarquables
\[
  \begin{aligned}
  \text{(43)} \qquad \mathrm{Gl}(2, \mathbb{Z})/\pm 1
  &\xrightarrow{\ \sim\ } \Pi_{0,3}^{\mathbb{D}_3, \tau}
  \quad (\xleftarrow{\ \sim\ } \Pi_{0,3}^\tau / \langle \rho_1^2, \rho_\infty^3 \rangle) \\
  \text{(44)} \qquad \mathrm{Sl}(2, \mathbb{Z})/\pm 1
  &\xrightarrow{\ \sim\ } \Pi_{0,3}^{\mathbb{D}_3}
  \quad (\xrightarrow{\ \sim\ } \Pi_{0,3} / \langle \rho_1^2, \rho_\infty^3 \rangle)
  \end{aligned}
\]
\note{le numéro (40) renvoie aux isomorphismes de la p. 71, qui n'y portent pas
de numéro lisible ; le (2) renvoie à une page antérieure à ce lot. La dernière
flèche de (44) va vers la droite sur la page, celle de (43) vers la gauche.}

\page{73}
permettant d'interpréter les éléments remarquables
$\rho, \sigma_i, \tau_i, \tilde{\sigma}_i, \struck{\ill{}}, \varepsilon_i, h_i$
de $\Pi_{03}^{\mathbb{D}_3, \tau}$, comme des éléments de
$\mathrm{Gl}(2, \mathbb{Z})/\pm 1$ --- que les \uncertain{eux} \uncertain{aussi}
\uncertain{remontent} de façon canonique en des éléments de
$\mathrm{Gl}(2, \mathbb{Z})$, donnant lieu : une formulaire des plus
sympathiques, que j'ai développé et utilisé intensivement dans mes notes de
l'an dernier (1981). Si on veut généraliser ceci dans le
\struck{\ill{}} sens que les iso.\ (40), au cas des groupes
\struck{\ill{}} $\Pi_{0n}$ \add{et variantes} associés aux polygones
réguliers \uncertain{(\ill{} $n \geq 3$)} il me semble qu'il faudra
\uncertain{trouver} une \uncertain{interprétation} des groupes
\uncertain{analogues} : $\mathrm{Sl}(2, \mathbb{Z})/\pm 1$ resp.\
$\mathrm{Gl}(2, \mathbb{Z})/\pm 1$, comme sous-groupes discrets
\add{remarquables} des groupes $\mathrm{Gl}(2, \mathbb{R})$…

(5) Dans l'interprétation « cartographique », laissons --- \uncertain{plutôt}
(40) s'explicitant comme deux équivalences de catégories, entre la cat.\
isotopique des cartes « triangulées » \add{$X$} (i.e.\ toutes \struck{\ill{}}
les faces sont d'ordre divisant trois, i.e.\ sont des triangles),
\struck{éventuelles} \add{et celle des} \struck{orientées} cartes triangulées
pondérées \add{$X'$} \emph{munies} d'une

\page{74}\note{p. 13 de l'auteur.}
opération de $\mathfrak{S}_3$ sur la carte pondérée sous-jacent, « échangeant
les pondérations », suivant l'action de $\mathfrak{S}_3$ sur l'ensemble
pondérant $\{0, 1, \infty\}$ --- équivalence tant dans la théorie orientée, que
dans la théorie non orientée. On peut expliciter cette équivalence,
\add{dans le cas orienté pour fixer les idées,} en disant que dans
\uncertain{un} sens : $X'$ on fait correspondre $X = X'/\mathfrak{S}_3$, avec la
« triangulation » déduite de la triangulation (pondérée) de $X'$ en passant au
quotient par $\mathfrak{S}_3$ --- on \struck{\uncertain{trouve}}
\add{trouve une} \ill{} qui aux \uncertain{centres} \ill{} une ramification,
\struck{\ill{}} plus, aux centres des faces (avec ramification
\emph{divis.\ à 3}), des \uncertain{centres} d'arêtes (ramification divis.\
à 2), et \uncertain{aux} sommets \ill{} de la triangulation pondérée
\ill{} ordinaire, \struck{\ill{} et dans les \ill{}} on considère la subdivision
barycentrique, \struck{\ill{}} \ill{} des \ill{} de celle-ci, on obtient de
l'un des \uncertain{sommets} des triangles de
\marginal{l'orientation \uncertain{servi} \uncertain{des} \ill{} :
\uncertain{rejoindre} \ill{} \uncertain{ramification} \ill{} des carte
\ill{} \ill{} \ill{} \uncertain{orientables} : la \uncertain{multiplicité} est
\uncertain{orientée} \ill{} \ill{} triangulaire \ill{} \ill{} \ill{}}
\note{note marginale écrite en diagonale dans la marge gauche, en grande partie
illisible.}
\ill{} \uncertain{principal} \struck{d'ordre} du groupe $\mathfrak{S}_3$,
défini \ill{} par la « couleur » sur chaque \ill{} défini \uncertain{entièrement}
par \uncertain{son} \uncertain{voisin} de la des pieces (qui en
\uncertain{composent}, \uncertain{comme} un \uncertain{voisin} des
\add{aux} \uncertain{principe} de la face qui le contient, \uncertain{par}
donné une pondération de ladite face…), et par le passage d'un drapeau
\add{aux} drapeaux adjacents\add{\struck{$\tau_0(r), \tau_1(r), \tau_\infty(r)$}}
(on trouve ainsi un \struck{\ill{} $\tau_0, \tau_1, \tau_\infty$}
revêtement ramifié principal \add{$X'$} de $X$, de groupe $\mathfrak{S}_3$,
dont lequel l'image inverse de la triangulation \uncertain{de départ}

\page{75}
est naturellement pondérée, c'est $X'$. Si la
\add{« triangulaire »} structure triangulaire $X$ est \uncertain{orientée}
(par des moyens d'orientés plus) dans la variante \uncertain{non}
\ill{} $X'$ n'est \uncertain{ramifiée} --- plus que les sommets de façon
générale \struck{\ill{}} \ill{} il y a ramification en un sommet
\add{de $X$} (ssi celui-ci est d'ordre impair. Comme toutes les ordres des
sommets \uncertain{sont} pairs, \add{les} on plus du fait que toutes les
faces sont des triangles et qu'il y a pas d'arêtes
\add{\uncertain{multiples}} \add{\uncertain{conditions}} plus
(\struck{toutes}) nécessaires pour que la structure \add{de carte}
triangulaire sur $X$ soit pondérable),
dans le revêtement $X'$ de $X$ est non ramifié. Les orientations,
correspondant 1--1 aux pondérations de la structure triangulaire sur $X$,
s'y \uncertain{trouvent}.

On peut \struck{\ill{}} \uncertain{décrire} encore un peu plus la
situation, dans le cas orienté que nous intéressons ici, en disant qu'une
condition nécessaire de pondérabilité est l'existence d'un « coloriage »
en deux « couleurs » $+$ et $-$, des faces par deux, des faces adjacentes
ayant des couleurs différentes. Les coloriages correspondent aux sections du
fibré \add{$X'/\mathfrak{S}_3^{+}$} associé au fibré de groupe
$\mathfrak{S}_3$, par $\mathfrak{S}_3 \xrightarrow{\ \mathrm{sg}\ } \pm 1$,
\uncertain{dont} \uncertain{un} \uncertain{tordu}\ill{}
\note{la phrase se poursuit sur la page suivante.}

\page{76}\note{p. 14 de l'auteur.}
au groupe alterné de $\mathfrak{S}_3$ : $\mathfrak{S}_3^{+} \simeq \mathbb{Z}/3\mathbb{Z}$.
Quand \struck{toutes} les conditions précédentes sont satisfaites, y compris la
donnée d'un \add{tel} coloriage des \struck{\ill{}}, on trouve un revêtement
principal, \uncertain{canonique}
\struck{$X'^{+} \subset X$ \ill{} de $X$}, de $X$
\[
  \text{(*)} \qquad X'^{+} \subset X'
\]
de groupe $\mathbb{Z}/3\mathbb{Z}$ --- dont la classe d'isomorphie est un élément
\[
  \text{(43)} \qquad \mathrm{cl}(X) \in H^1(X, \mathbb{Z}/3\mathbb{Z})
\]
qu'on peut considérer comme l'obstruction : l'existence d'une pondération, pour
la carte triangulaire $X$ : la nullité est nécessaire et suffisante pour cette
existence. De plus, les pondérations pour lesquelles les faces $+$ sont les faces
directes (pour l'orientation donnée) sont en corr.\ 1--1 avec les sections du
\struck{fibré} torseur principal précédent de groupe $\mathbb{Z}/3\mathbb{Z}$.
\note{ce numéro (43) fait double emploi avec le (43) de la p. 72. Le reste de la page est blanc, barré d'un trait oblique.}

\page{77}\note{p. 15 de l'auteur.}
(6) Soit
\[
  \text{(44)} \qquad P = (S, A, R \subset S \times A)
\]
un polygone combinatoire fini d'ordre $n \geq 3$,
\[
  \text{(45)} \qquad \mathbb{D}_P = \mathbb{D}
\]
le groupe de ses automorphismes, $\mathbb{D}_P^{+} = \mathbb{D}^{+}$ le sous-groupe
d'indice 2 des automorphismes conservant l'orientation, ou rotations,
\[
  \text{(46)} \qquad 1 \longrightarrow \mathbb{D}^{+} \longrightarrow \mathbb{D}
  \longrightarrow \{\pm 1\} \longrightarrow 1 .
\]
\note{le numéro (46) se lit (45) sur la page, qui porte ainsi deux (45).}
On identifie l'ensemble $\underline{\omega}_P = \underline{\omega}$ des
orientations de $P$ : l'ens.\ des deux générateurs privilégiés de $\mathbb{D}^{+}$,
\struck{\uncertain{conservant} \uncertain{l'orientation}}
\[
  \text{(47)} \qquad \underset{\{\omega, \omega'\}}{\underline{\omega}}
  \subset \mathbb{D}^{+}
\]
de sorte que pour $\omega \in \underline{\omega}$, et pour tout $s \in S$,
$a \in A$, $\omega(s)$ est un sommet \emph{adjacent} : $s$, $\omega a$ une arête
adjacente : $a$. Il existe un plongement, unique : isom.\ unique
\struck{\ill{}} près, de $P$ dans la sphère complexe $\Sigma = \Sigma_P$, de
telle façon que les opérations de $\mathbb{D}$ se prolongent en
\struck{\uncertain{anti}}automorphismes \struck{\ill{}} complexes de $\Sigma_P$ :
\[
  \text{(48)} \qquad S = S_P \hookrightarrow \Sigma = \Sigma_P
\]
--- l'unicité est \uncertain{déjà} garantie en exigeant le prolongement des
opérations de $\mathbb{D}^{+}$, qui implique celle des opérations de $\mathbb{D}$
tout entier. Alors le polygone $P$ est réalisé canoniquement par des points de
$\Sigma$ et des arcs de \add{grand} cercle de $\Sigma$, \struck{\ill{}}
\struck{derniers} \add{tous} placés sur un même (grand) cercle,
\marginal{NB Cela revient : la donnée d'un \ill{} \uncertain{corps} \ill{} des
\uncertain{groupes} \ill{} \ill{} \ill{} d'un \ill{} \ill{} \ill{}
\uncertain{orienté} \ill{} \ill{} \ill{} \ill{} \ill{} --- La démonstration se
fait en \ill{} \ill{} \ill{} puis \ill{} $\mathbb{C}^{*}$ \ill{}
$0, \infty$ \ill{} $\mathbb{P}^1_{\mathbb{C}}$ \ill{} automorphismes
\ill{} \ill{} \ill{} $\Sigma^{0} = \mathbb{P}^1_{\mathbb{C}}$ \ill{}
$\mathbb{D}_P^{+}$ \ill{} \ill{} \ill{} (48) \ill{} \ill{} $\mathbb{D}_P^{+}$
\ill{} points \ill{} $R_{\infty}$ \ill{} \ill{} $P_\omega \mapsto \infty$ ?
\ill{} \ill{}}
\note{la longue note marginale, écrite en diagonale le long de la marge gauche
et descendant jusqu'au bas de la page, est en très grande partie illisible ;
on n'en donne que les fragments sûrs.}
\drawing{En bas à gauche, une sphère au crayon : l'équateur tracé en ellipse, avec des points marqués dessus (le polygone $P$ plongé sur un grand cercle).}

\page{78}
dont ils forment une subdivision cellulaire --- ce cercle définit une structure
réelle canonique sur la droite projective complexe $\Sigma$,
\add{il est donc (cf.\ \uncertain{dimensions} \uncertain{constantes} \uncertain{près})}
\struck{\ill{}} de $\Sigma_{\mathbb{R}}$, notée $\Sigma(\mathbb{R})$.
\struck{\ill{} $\Sigma_{\mathbb{R}}$} De plus, il y a une métrique riemannienne
homogène unique sur $\Sigma$ invariante par $\mathbb{D}$ (c'est pour elle que
$\Sigma(\mathbb{R})$ est un \emph{grand} cercle).
$\Sigma_{\mathbb{R}}(\mathbb{R})$ partage $\Sigma$ en deux hémisphères,
\struck{\ill{}} en correspondance bijective canonique avec $\underline{\omega}$
par la correspondance de Stokes
\[
  \text{(49)} \qquad \pi_0(\Sigma \smallsetminus \Sigma_{\mathbb{R}}) \simeq
  \underline{\omega} .
\]
On a
\[
  \Sigma^{\mathbb{D}} = \Sigma^{\mathbb{D}^+} = \Sigma^{\underline{\omega}}
  = \Sigma^{\{\omega\}} \subset \Sigma \smallsetminus \Sigma_{\mathbb{R}}
  \qquad (\omega \in \underline{\omega})
\]
et $\operatorname{card} \Sigma^{\mathbb{D}} = 2$, de façon précise, chacun des
hémisphères \uncertain{ouverts} contient un pt et un seul de $\Sigma^{\mathbb{D}}$,
d'où
\[
  \text{(50)} \qquad \Sigma^{\mathbb{D}} \simeq \underline{\omega} ,
\]
on désigne par $P_\omega$ le point de $\Sigma^{\mathbb{D}}$ qui correspond :
$\omega \in \underline{\omega}$. Il est caractérisé par le fait que l'opération de
$\omega$ dans l'espace tangent : $\Sigma$ en $P_\omega$ est une rotation d'angle
$+2\pi/n$ (alors qu'en $P_{\omega'}$ elle est d'angle $-2\pi/n$). Pour $s \in S$,
on désigne par $R_s$ le point correspondant de $\Sigma_{\mathbb{R}}$, pour
$a \in A$ on désigne par $Q_a$ le centre de l'arête $a$.

\page{79}\note{p. 16 de l'auteur.}
\drawing{En haut à gauche, une sphère : $P_\omega$ au pôle nord, $P_{\omega'}$ au pôle sud ; sur l'équateur, alternativement des points $R_s$, $R_{\omega s}$, $R_{\omega' s}$ (gros points) et des centres d'arêtes $Q_a$, $Q_{\omega a}$, $Q_{\omega' a}$ (croix) ; des demi-méridiens en pointillé $\lambda_a$, $\lambda_{\omega a}$, $\lambda_{\omega' a}$ joignant $P_{\omega'}$ à $P_\omega$ en passant par les $Q$ ; une flèche $\tau_\omega$ \uncertain{autour} de l'équateur. Sous la figure : « $(s, a) = r$ d'orientation $\omega$ ».}

Je m'intéresse aux groupoïdes fondamentaux
\[
  \Pi_1(\underbrace{\Sigma_P \smallsetminus S_P}_{\Sigma_P^{*} \text{ ou } \Sigma^{*}}) ,
\]
et je vais écrire d'abord le groupe fondamental en un $P_\omega$, correspondant
\uncertain{en} \uncertain{donnée} de $\omega \in \underline{\omega}$, \uncertain{qu'on}
\uncertain{note}
\[
  \text{(51)} \qquad
  \underset{\text{ou } \Pi_P \text{ ou } \Pi}{\Pi_{P,\omega}}
  = \Pi_1(\Sigma_P^{*}, P_\omega) .
\]
On considère aussi l'opération antiholomorphe unique qui fixe les pts de
$\Sigma_{\mathbb{R}}$, et qui échange $P_\omega$ et $P_{\omega'}$, c'est la
conjugaison complexe pour la structure réelle canonique de $\Sigma$, on la
note $\tau$.

Ainsi, le groupe
\[
  \text{(52)} \qquad \mathbb{D}_P \times \{1, \tau\} \qquad \text{d'ordre } 4n
\]
\note{le numéro se lit (52) ou (51) ; le suivant est (53).}
opère sur $\Sigma$, en respectant la « carte cyclotomique » sur $\Sigma$, le
sous-groupe qui conserve l'orientation \uncertain{étant} $\mathbb{D}_P$, alors
qu'il induit l'identité sur $S$ (ou sur $P$) \uncertain{est} $\{1, \tau\}$.
On va aussi regarder les groupoïdes fondamentaux mixtes de $\Sigma^{*}$, pour
l'action sur $\Sigma^{*}$ de $\{1, \tau\}$, de $\mathbb{D}_P$, ou de
$\mathbb{D}_P \times \{1, \tau\}$. \add{Les} \struck{Les groupes fondamentaux}
\add{\uncertain{posons}}
\[
  \text{(53)} \qquad
  \left\{
  \begin{aligned}
    \Pi_{P,\omega}^{\tau} &= \Pi_1(\Sigma^{*}, \{1, \tau\}; P_\omega) \\
    \Pi_{P,\omega}^{\mathbb{D}} &= \Pi_1(\Sigma^{*}, \mathbb{D}; P_\omega) \\
    \Pi_{P,\omega}^{\mathbb{D},\tau} &= \Pi_1(\Sigma^{*}, \mathbb{D} \times \{1, \tau\}; P_\omega)
  \end{aligned}
  \right.
\]
\note{dans la deuxième ligne, l'exposant porte une marque surchargée, peut-être $\mathbb{D}^{*}$.}
qu'on écrira aussi $\Pi^\tau, \Pi^{\mathbb{D}}, \Pi^{\mathbb{D},\tau}$ quand $P$
est sous-entendu. On a donc \struck{des} une des groupes d'extensions de

\page{80}
groupes fondamentaux (54) :

\begin{tikzcd}[column sep=small, row sep=small]
 & \Pi \arrow[dl, hook] \arrow[dr, hook] & \\
 \Pi^{\mathbb{D}} \arrow[dr, hook] & & \Pi^{\tau} \arrow[dl, hook] \\
 & \Pi^{\mathbb{D},\tau} &
\end{tikzcd}

\note{dans le losange (54), la flèche de $\Pi$ vers $\Pi^{\tau}$ porte le crochet d'inclusion du côté de $\Pi$ mais sa pointe n'est pas visible ; on la transcrit comme les trois autres.}

Pour exprimer commodément des générateurs remarquables de ces groupes, on
introduit des chemins \add{\struck{\ill{}}} \add{les} \struck{\ill{}}
indexés par $a \in A$, i.e.\ les \emph{arêtes} du polygone)
\[
  \text{(55)} \qquad \lambda_a : P_{\omega'} \longrightarrow P_\omega
  \quad \text{arc de demi-grand cercle passant par } Q_a
\]
\drawing{Dans la marge gauche, une sphère : $P_\omega$ en haut, $P_{\omega'}$ en bas, deux demi-méridiens en pointillé $\lambda$ passant par $Q_{\omega' a}$ et $Q_a$ sur l'équateur, avec $R_s$ entre eux.}

Pour tout $s \in S$, soit
\[
  \left\{
  \begin{array}{l}
    a(s, \omega) \text{ ou } a(s) \in A \text{ tel que } (s, a) \text{ soit un repère} \\
    \text{d'orientation } \omega, \text{ i.e.\ } s = \mathrm{or}_\omega(a)
  \end{array}
  \right.
\]
\note{après « i.e. », une première formule est biffée et illisible ; $s = \mathrm{or}_\omega(a)$ est ajouté au-dessus.}
on pose
\[
  \text{(56)} \qquad \rho_s = \lambda_a \lambda_{\omega' a}^{-1} ,
  \qquad \text{où } a = a(s, \omega) ,\ \omega' = \omega^{-1}
\]
\[
  \text{donc} \quad \omega'(a) = a(\omega'(s), \omega) ,
\]
\[
  \underset{\Pi_P = \Pi_1(\Sigma_P^{*}, P_\omega)}{\cap}
\]
\note{le signe $\cap$ (« appartient ») est placé sous $\rho_s$, avec en dessous $\Pi_P = \Pi_1(\Sigma_P^{*}, P_\omega)$.}
On trouve alors, pour tout $s \in S$
\[
  \text{(57)} \qquad \rho_{\omega^{n-1}s}\, \rho_{\omega^{n-2}s} \cdots
  \rho_{\omega s}\, \rho_s = 1 .
\]
Toutes ces \add{$n$} relations sont équivalentes entre elles. Si on choisit un
sommet origine $s_0$ --- ce qui, une fois choisi $\omega$, fait un \emph{repère}
\struck{de} $r = (s_0, a(s_0, \omega))$ de $P$, donc aussi un repère de la carte
cyclotomique associée), on pose aussi
\[
  \text{(58)} \qquad \rho_i = \rho_{\omega^i s_0}
  \qquad (i \in \mathbb{Z}/n\mathbb{Z})
\]
--- les indices sont entiers le plus souvent $0, 1, \ldots, n-1$ ---
de sorte que (57), prenant la forme…
\note{la phrase se poursuit au-delà de ce lot, à la page 81.}
\end{document}
